\chapter{Final Design}

\section{Introduction to Final Design}

Following the evaluation of design alternatives and the establishment of governing design criteria, the wastewater conveyance system for the Colebrooke Subdivision is finalized as an integrated lift station and force main system. The selected configuration consists of a duplex submersible pump station that conveys wastewater through a pressurized force main to an existing municipal gravity sewer manhole located along Lewis Street. The system layout, including the pump station location, force main alignment, and downstream connection, is defined based on the subdivision survey and utility plan \cite{ProjectD_Colebrooke_2020}.

The force main extends $77.48~\text{ft}$ from the pump station to the receiving municipal manhole. Wastewater is collected from the subdivision gravity sewer network and delivered to the wet well, where it is pumped to the downstream system. The governing hydraulic input for the design is a peak wastewater inflow of $12.46~\text{gpm}$. The vertical difference between the wet well operating level ($299.00~\text{ft}$) and the receiving sewer invert ($309.28~\text{ft}$) produces a static head of $10.28~\text{ft}$, which forms the primary component of the system head requirements.

Based on the hydraulic analysis and system curve evaluation, the pump operating condition is established at $135.5~\text{gpm}$ and $12.08~\text{ft}$ of head. These parameters define the required performance of the pumping system and ensure that wastewater is reliably conveyed under expected operating conditions. This chapter presents the detailed hydraulic calculations, force main design, pump selection, and wet well sizing that form the basis of the final system design.


